\documentclass[12pt]{article}
\usepackage[utf8]{inputenc}
\usepackage[T1]{fontenc}
\usepackage{hyperref}
\title{Lolopo, Southern Language Resources for K-12 Curriculum Development}
\author{Compiled by Claude Sonnet 5.5 (claude-sonnet-5-5) for Andr\'as Kornai}
\date{}
\begin{document}
\maketitle

\section{Introduction}
Southern Lolopo (ISO 639-3 \texttt{ysp}, Glottocode sout3129) is a Loloish (Sino-Tibetan) language of central Yunnan, China. It is part of the Lolopo group (ycl), which Yi people speak. Estimates of the speaker population conflict. One Loloish table gives 190,000, and Joshua Project gives 250,000 for the Southern Luoluopo group. The search found very few resources, and none for use in schools.

\section{Primary Educational Resources}
No primers, textbooks or school programmes for ysp were found.

\section{Linguistic and Reference Materials}
\begin{itemize}
\item Wikipedia, Lolopo language, an overview of the Lolopo group: \url{https://en.wikipedia.org/wiki/Lolopo_language}
\item Wikipedia, Loloish languages, with a table of the languages and their speaker numbers: \url{https://en.wikipedia.org/wiki/Loloish_languages}
\end{itemize}

\section{Religious and Community Resources}
\begin{itemize}
\item Joshua Project, Southern Luoluopo in China. Bible translation is listed as needed. The only resource listed is GRN audio teaching: \url{https://joshuaproject.net/people_groups/18569}
\end{itemize}

\section{Digital and Audio Resources}
\begin{itemize}
\item Global Recordings Network page for ysp. It states that no recordings are currently available: \url{https://globalrecordings.net/en/language/ysp}
\item Global Recordings Network, ``Yi: Lolo in Mojiang''. This is a related variety linked to ysp, with audio Bible stories. I did not verify that it is the same as ysp: \url{https://globalrecordings.net/en/language/6077}
\end{itemize}

\section{Curriculum Development Context}
No published teaching materials, orthography documents or dictionaries for ysp were found, and this search did not look at the Chinese-language literature. A curriculum would have to start from new documentation. Standard Yi (Nuosu) materials may be a model, but they describe a different language.
\end{document}